\section{Revisión de orden $h$: fronteras estructurales, complejidad y mezcla}
\label{sec:cfrd}

La dinámica de revisión mueve a lo sumo $h$ AMR por paso. Esta sección
responde tres preguntas sobre ese número: cuándo $h=2$ basta y cuándo no,
qué cuesta encontrar la revisión mínima cuando no basta, y qué determina la
velocidad —no solo la alcanzabilidad— de la revisión estocástica
\cite{mayorga2026compendio}.

\subsection{Por qué hace falta orden dos, y por qué la factibilidad ya es dura}

\begin{lema}[Barrera unilateral y escape bilateral]
\label{le:barrera-unilateral}
Sean dos cargas con demandas de capacidad $\bar c_\ell,\bar c_{\ell'}$
cubiertas por dos AMR $i\in\ell$, $j\in\ell'$ con capacidades $c_i,c_j$, de
modo que las coberturas actuales son $\sigma_\ell$ y $\sigma_{\ell'}$. Si
$\sigma_\ell-c_i<\bar c_\ell$, $\sigma_{\ell'}-c_j<\bar c_{\ell'}$,
$\sigma_\ell-c_i+c_j\ge\bar c_\ell$, $\sigma_{\ell'}-c_j+c_i\ge\bar c_{\ell'}$
y el intercambio reduce el coste, entonces ninguno de los dos movimientos que
lo componen mejora por separado —cada uno deja una carga descubierta— y
aplicarlos a la vez mantiene la cobertura y reduce el coste.
\end{lema}

\begin{proof}
Las dos primeras desigualdades dicen que retirar $i$ de $\ell$ o $j$ de
$\ell'$ crea déficit y empeora el orden lexicográfico (déficit antes que
coste); las dos siguientes, que tras el intercambio ambas quedan cubiertas;
la última, que el coste baja.
\end{proof}

\begin{observacion}[La escalera de órdenes empieza aquí]
\label{obs:escalera-ordenes}
\Cref{le:barrera-unilateral} es el caso $h=2$ de lo que
\cref{th:compatibilidad-completa} lleva a $h$ arbitrario: un intercambio que
ningún movimiento individual realiza. Y hay dos cotas de complejidad
distintas que conviene no confundir. Decidir si existe una asignación
factible con dos cargas y capacidades aditivas es ya NP-duro, por reducción
de partición: con dos cargas de demanda $\sigma/2$ y AMR de capacidades
$c_1,\dots,c_n$ con $\sum c_i=\sigma$, hay asignación factible si y solo si
los $c_i$ admiten una partición en dos subconjuntos de suma $\sigma/2$; una
carga sola no bastaría, porque asignarle toda la flota siempre cubre. Y
encontrar la reparación de orden mínimo es NP-completo y W[2]-completo por
\cref{th:mohr}. La primera dureza es de la factibilidad; la segunda, de la
optimalidad del recurso; la construcción $C_{2h}$ muestra que la revisión
puede ser difícil aun cuando la factibilidad sea trivial.
\end{observacion}

\begin{proposicion}[Composición de propuestas disjuntas]
\label{pr:composicion-disjunta}
Propuestas que no comparten AMR ni cargas, de origen ni de destino, suman sus
variaciones de déficit y de coste; si todas mejoran en orden lexicográfico, el
lote también. Y si todos los candidatos se examinan y cada \emph{commit}
reduce el mérito lexicográfico, la cadena termina en un mínimo $h$-local en a
lo sumo $|A|-1$ cambios sobre el espacio finito $A$.
\end{proposicion}

\begin{proof}
Déficit y coste suman sobre cargas y AMR; con soportes disjuntos cada
propuesta altera sumandos distintos. La terminación es la de
\cref{th:potencial}: ningún perfil se repite en un orden estricto sobre un
conjunto finito.
\end{proof}

\subsection{Cuándo la revisión por pares no basta}

Una objeción natural a los ejemplos de orden alto es que reflejen una matriz
de compatibilidad artificialmente dispersa. La construcción siguiente la
elimina.

\begin{teorema}[Compatibilidad completa no implica reconfigurabilidad por pares]
\label{th:compatibilidad-completa}
Para todo $h\ge3$ existe un problema de asignación con $h$ tareas,
compatibilidad física individual completa y restricciones aditivas de
capacidades tal que la fibra factible contiene exactamente dos asignaciones de
los AMR pivote y la distancia de revisión entre ellas es $h$.
\end{teorema}

\begin{proof}
Cada tarea tiene un AMR ancla fijo y hay $h$ AMR pivote. Con $h$ capacidades
binarias, el pivote $j$ aporta $p_j=e_j+e_{j-1}$ (índices módulo $h$) y la
tarea $k$, contada su ancla, necesita el residual $e_k$. Todos los pivotes son
físicamente compatibles con todas las tareas, pero por capacidad la tarea $k$
solo puede cubrirse con los pivotes $k$ o $k+1$: el grafo bipartito
pivote--tarea es el ciclo par $C_{2h}$, que tiene exactamente dos
emparejamientos perfectos, y difieren en los $h$ pivotes. No hay transición
factible de orden menor que $h$, y la rotación completa sí es factible. La
familia se construyó para $h=3,\dots,10$: en todos los casos dos
emparejamientos y orden mínimo exactamente $h$.
\end{proof}

\begin{observacion}[La dificultad viene del acoplamiento, no de la dispersión]
\label{obs:fragmentacion-aleatoria}
Con seis AMR, tres tareas, dos AMR por tarea, tres capacidades aditivas y
compatibilidad completa, la fragmentación del grafo de revisión por pares
aparece sola al apretar la demanda: con holgura $0{,}6$ en dos de mil
mundos, con $0{,}8$ en el $7{,}9\%$ y con $0{,}9$ en el $10{,}8\%$
\cite{mayorga2026compendio}. La conexión con \cref{th:saltos} es directa: el
orden de revisión necesario es una propiedad de la fibra de capacidades, y
\cref{th:compatibilidad-completa} dice que puede ser arbitrariamente alto sin
que ningún AMR sea incompatible con nada.
\end{observacion}

\begin{observacion}[Revisión conjunta medida contra el óptimo enumerado]
\label{obs:revision-conjunta-r11}
Con ocho AMR, dos cargas y cuatro miembros por carga hay $\binom84=70$
particiones, luego el óptimo es enumerable. En cien instancias heterogéneas
con dos recursos aditivos y eventos que cambian costes, reducen capacidad o
suben demanda, doce quedan globalmente inviables y se conservan como tales;
en las ochenta y ocho factibles, el intercambio de un par encuentra el óptimo
en $75$, hasta dos pares en $86$, y la ampliación adaptativa del orden en
$88$: trece óptimos adicionales —siete reparaciones que el par no encuentra y
seis factibles pero subóptimas—, McNemar exacto $p=2^{-12}=2{,}4\times10^{-4}$,
que sobrevive a Holm \cite{mrob2026r11}. El contraejemplo reproducible es de
manual: la capacidad de un AMR cae de $(8,8)$ a $(2,2)$, los pares se detienen
en $48{,}88$ y una revisión conjunta de cuatro llega a $38{,}19$, un $21{,}9\%$
mejor. Es \cref{th:compatibilidad-completa} con costes: converger bien dentro
de un vecindario insuficiente no resuelve el problema.

Y el límite, igual de medido. En cien instancias homogéneas los pares ya dan
$100/100$: ampliar no añade calidad. Y el procedimiento adaptativo que
certifica que no queda mejora inspecciona en mediana $85$ consultas al
evaluador frente a $70$ de enumerar directamente: no ha demostrado evitar la
explosión combinatoria de \cref{th:mohr}, y una búsqueda genérica de
vecindario variable con los mismos soportes da la misma respuesta en las
doscientas instancias. La ganancia se atribuye al soporte de revisión, no al
vocabulario; el aporte pendiente es descubrir el bloque necesario con menos
consultas que enumerarlo, que es \cref{obs:dicotomia} vuelta problema
algorítmico.
\end{observacion}

\begin{observacion}[La frontera positiva: convexidad discreta]
\label{obs:m-convexo}
El caso contrario también está caracterizado. Si el dominio de asignaciones es
M-convexo y el objetivo es una función M-convexa, un punto es óptimo global si
y solo si ningún intercambio elemental $x\to x-e_i+e_j$ lo mejora
\cite{murota2003discrete}; si cada intercambio elemental se realiza con una
revisión de a lo sumo dos AMR, todo terminal por pares es óptimo global y los
términos topológico y dinámico de \cref{th:descomposicion} son cero. La
necesidad de orden alto es, por tanto, un indicador de pérdida de estructura
de intercambio —que las cuotas y capacidades acopladas destruyen— y no un
objetivo en sí.
\end{observacion}

\subsection{Cuánto cuesta la reparación mínima}

\begin{proposicion}[Reparación por ranuras con recurso mínimo]
\label{pr:reparacion-ranuras}
Si cada requisito es una ranura unitaria y cada AMR ocupa a lo sumo una, tras
un fallo la reparación con el mínimo número de AMR reasignados es un camino
aumentante más corto desde la ranura vacía hasta un AMR libre en el grafo
residual alternante, y se encuentra por búsqueda en anchura en
$O(|V|+|E|)$.
\end{proposicion}

\begin{proof}
Un camino alternante de longitud $h$ induce exactamente $h$ reasignaciones y
conserva la ocupación de las ranuras intermedias. Recíprocamente, la
diferencia simétrica entre el emparejamiento incompleto y cualquier reparación
contiene un camino aumentante desde la ranura vacía hasta un AMR antes libre,
cuyo número de AMR no excede el recurso de la reparación.
\end{proof}

\begin{teorema}[La reparación heterogénea mínima es NP-completa]
\label{th:mohr}
Sea el problema: dada una tarea con requisitos binarios $b_m=1$, AMR libres
con capacidades $a_{im}\in\{0,1\}$ y un entero $h$, decidir si existe una
coalición de a lo sumo $h$ AMR con $\sum_{i\in\Ccal}a_{im}\ge1$ para todo $m$.
Es NP-completo incluso con una sola tarea, sin batería, sin costes y sin
restricciones de comunicación; parametrizado por $h$ es W[2]-completo, y su
versión de minimización hereda la aproximación voraz de factor
$H_m\le1+\ln m$.
\end{teorema}

\begin{proof}
Pertenencia a NP: dada la coalición se verifican las desigualdades en tiempo
polinomial. Dureza: de recubrimiento de conjuntos \cite{karp1972reducibility},
cada elemento del universo es un requisito $m$ y cada conjunto $S_i$ un AMR
con $a_{im}=1$ si y solo si $m\in S_i$; hay recubrimiento de cardinalidad
$\le h$ si y solo si hay coalición factible de cardinalidad $\le h$. La
reducción conserva el parámetro, y recubrimiento parametrizado por $h$ es
W[2]-completo \cite{downey2013fundamentals}.
\end{proof}

\begin{observacion}[La dicotomía operativa]
\label{obs:dicotomia}
Reparación por ranuras o dominio de intercambio convexo: reparación exacta
local, polinomial, por pares. Capacidades acopladas generales: reparación de
orden mínimo NP-dura y no tratable en $f(h)\,n^{O(1)}$ salvo colapso
FPT$=$W[2]. Es la justificación algorítmica de usar $h$ como presupuesto de
escalada y no como promesa de optimalidad: un bloque óptimo pequeño no hace el
problema fácil. La caracterización matricial exhaustiva de cuándo basta $h=2$
queda abierta.
\end{observacion}

\subsection{Del continuo al atómico con probabilidad}

\begin{proposicion}[Estrechamiento probabilístico continuo--atómico]
\label{pr:hoeffding}
Sean decisiones atómicas independientes $S_i$ muestreadas de un perfil
continuo, con carga $Y_z=\sum_{i=1}^{N}a_{iz}S_{iz}$ sobre cada recurso $z$ y
$0\le a_{iz}S_{iz}\le\bar a_{iz}$. Si
\begin{equation}
  \label{eq:hoeffding}
  \mathbb E[Y_z]\le C_z-\varepsilon_z,\qquad
  \varepsilon_z=\sqrt{\tfrac12\sum_i\bar a_{iz}^2\,\log\tfrac{|Z|}{\delta}},
\end{equation}
entonces $\Pr\{Y_z\le C_z\ \forall z\}\ge1-\delta$.
\end{proposicion}

\begin{proof}
Hoeffding da $\Pr\{Y_z-\mathbb EY_z\ge\varepsilon_z\}\le\exp(-2\varepsilon_z^2/\sum_i\bar a_{iz}^2)=\delta/|Z|$;
una cota de unión sobre los $|Z|$ recursos concluye. Con doscientas
decisiones, cuatro recursos y $\delta=0{,}05$ la violación observada por
recurso es $4\times10^{-4}$ frente a la garantía de $1{,}25\times10^{-2}$.
\end{proof}

\begin{observacion}[Es el margen de \cref{co:tensado}, en la capa discreta]
\label{obs:hoeffding-uso}
El estrechamiento de \cref{co:tensado} convierte covarianza de estado en
margen geométrico; \eqref{eq:hoeffding} convierte la aleatoriedad del cierre
atómico en margen de capacidad. Ambos pagan lo mismo —recurso no usado— por
lo mismo: una garantía en probabilidad sobre lo que el modelo continuo no
resuelve. Y la cota es conservadora por un orden de magnitud, que es el precio
de no suponer nada sobre la distribución.
\end{observacion}

\subsection{Transferencias: cuándo un bloque beneficioso es aceptable}

Que el potencial suba no basta para agentes egoístas: la desviación de dos
agentes con cambios de utilidad $(2,-1)$ tiene superávit $1>0$ y el segundo la
rechaza sin compensación.

\begin{proposicion}[Implementabilidad con presupuesto equilibrado]
\label{pr:implementabilidad-tu}
Sea $\Ccal$ una coalición que ejecuta un bloque con cambios de utilidad
pretransferencia $d_i$ y ganancia de potencial $g=\sum_{i\in\Ccal}d_i$. Bajo
utilidad cuasilineal, existe un vector de transferencias $\theta$ con
$\sum_i\theta_i=0$ y $d_i+\theta_i>0$ para todo $i$ si y solo si $g>0$; y
con pesos $w_i>0$ el reparto que maximiza $\sum_iw_i\log s_i$ sujeto a
$\sum s_i=g$ es $s_i=w_ig/\sum_jw_j$.
\end{proposicion}

\begin{proof}
Necesidad: $g=\sum_i(d_i+\theta_i)>0$. Suficiencia: $s_i=w_ig/\sum_jw_j>0$
y $\theta_i=s_i-d_i$ suman cero. El reparto es la estacionariedad
$w_i/s_i=\lambda$ del lagrangiano.
\end{proof}

\begin{corolario}[Equilibrio fuerte de orden $h$ con transferencias]
\label{co:fuerte-tu}
Si toda coalición de tamaño $\le h$ puede redistribuir con presupuesto
equilibrado el cambio de potencial de su bloque, un estado no admite bloques
factibles de tamaño $\le h$ con ganancia positiva si y solo si ninguna
coalición de tamaño $\le h$ puede implementar una desviación factible que
mejore estrictamente a todos sus miembros mediante transferencias.
\end{corolario}

\begin{observacion}[Dos lecturas que no deben mezclarse]
\label{obs:flota-vs-estrategico}
En control de flota, el sistema impone un bloque globalmente beneficioso y
\cref{co:fuerte-tu} es innecesario. En la lectura estratégica es lo que
convierte el equilibrio fuerte de \cref{def:fuerte} en un concepto con
mecanismo: sin transferencias o sin regla eficiente de reparto la equivalencia
desaparece, y el reparto de \cref{sec:standby} es exactamente ese mecanismo.
\end{observacion}

\subsection{Guarda robusta de ganancia}

\begin{proposicion}[Aceptar solo lo que sobrevive al error]
\label{pr:guarda-ganancia}
Si la ganancia de un bloque se estima con estadísticas locales imperfectas,
$|\hat g-g|\le\bar\varepsilon_g$ con
$\bar\varepsilon_g=\sum_kL_k(\epsilon_k+\epsilon_k')+\sum_i\epsilon_i^{c}$
(constantes de Lipschitz de los valores por tarea y errores de estado y de
coste), y se aceptan solo bloques con $\hat g\ge\eta+\bar\varepsilon_g$,
entonces todo bloque aceptado tiene $g\ge\eta$ y el potencial real es monótono.
\end{proposicion}

\begin{proof}
$g\ge\hat g-\bar\varepsilon_g\ge\eta$; la cota de $\bar\varepsilon_g$ es la
desigualdad triangular término a término.
\end{proof}

\begin{observacion}[El estimador entrega la envolvente, la guarda la usa]
\label{obs:guarda-envolvente}
Es la misma lógica que \cref{pr:commit}: cualquier esquema de consenso o
estimación distribuida que entregue cotas certificadas $\epsilon_k(t)$ bajo
conectividad conjunta y retardo acotado —el filtro de \cref{sec:estimacion} o
el observador distribuido de \cref{sec:comunicacion}— se convierte en un
margen de aceptación. La convergencia del estimador es resultado previo; lo
que se aporta es usar su envolvente de error como guarda estratégica.
\end{observacion}

\subsection{Velocidad, no solo alcanzabilidad}

Dos soportes conectados pueden mezclar a velocidades muy distintas. Sobre un
grafo de revisión finito conectado $G_M=(\mathcal F,E_M)$ con potencial $\Psi$
y temperatura inversa $\beta$, considérese la cadena reversible con tasas
$q_{zz'}=a_{zz'}\exp\bigl(\beta(\Psi(z')-\Psi(z))/2\bigr)$, $a_{zz'}=a_{z'z}>0$,
cuya estacionaria es $\pi(z)\propto e^{\beta\Psi(z)}$.

\begin{teorema}[Disipación exacta de la energía libre]
\label{th:energia-libre}
Con $r_t(z)=\mu_t(z)/\pi(z)$ y $c_{zz'}=\pi(z)q_{zz'}=c_{z'z}$,
\begin{equation}
  \label{eq:energia-libre}
  \frac{\mathrm d}{\mathrm dt}D_{\mathrm{KL}}(\mu_t\,\|\,\pi)
  =-\frac12\sum_{z,z'}c_{zz'}\bigl[r_t(z')-r_t(z)\bigr]\bigl[\log r_t(z')-\log r_t(z)\bigr]\le0 .
\end{equation}
\end{teorema}

\begin{proof}
Derivar $D=\sum_z\mu_z\log r_z$ con conservación de masa da
$\dot D=\sum_z\dot\mu_z\log r_z$; la ecuación maestra y el balance detallado
simetrizan el flujo, y $(a-b)(\log a-\log b)\ge0$.
\end{proof}

\begin{teorema}[La brecha espectral crece con el soporte]
\label{th:gap-monotono}
Sean $M_1\subseteq M_2$ dos soportes de revisión sobre la misma fibra y el
mismo potencial, con los mismos relojes en las aristas comunes y relojes no
negativos en las adicionales. Con la forma de Dirichlet
$\mathcal E_M(f,f)=\tfrac12\sum c_{zz'}[f(z')-f(z)]^2$ y
$\gamma_M=\inf_{\mathrm{Var}_\pi f>0}\mathcal E_M(f,f)/\mathrm{Var}_\pi(f)$,
se tiene $\gamma_{M_2}\ge\gamma_{M_1}$; si $G_{M_1}$ es desconectado,
$\gamma_{M_1}=0$, y si $G_{M_2}$ es conectado, $\gamma_{M_2}>0$. Además
$\lVert\mu_t/\pi-1\rVert_{L^2(\pi)}\le e^{-\gamma_Mt}\lVert\mu_0/\pi-1\rVert_{L^2(\pi)}$
y $\lVert\mu_t-\pi\rVert_{TV}\le\tfrac12e^{-\gamma_Mt}\lVert\mu_0/\pi-1\rVert_{L^2(\pi)}$.
\end{teorema}

\begin{proof}
$\mathcal E_{M_2}=\mathcal E_{M_1}+\tfrac12\sum_{E_2\setminus E_1}c_{zz'}[f(z')-f(z)]^2\ge\mathcal E_{M_1}$
con la misma $\pi$, y el ínfimo del cociente de Rayleigh conserva la
desigualdad; en un grafo finito reversible, conectividad equivale a que cero
sea autovalor simple. Las cotas de mezcla son las estándar de cadenas
reversibles. Sobre treinta estados con potencial aleatorio, un camino y el
mismo camino con aristas añadidas: brecha $0{,}0014$ frente a $0{,}638$, y la
divergencia decrece en cada uno de los treinta instantes muestreados.
\end{proof}

\begin{proposicion}[Cota de cuatro mecanismos]
\label{pr:cuatro-mecanismos}
Sean $\Psi_{\mathrm{rel}}$ el máximo del potencial continuo,
$\Psi_{\mathrm{IP}}$ el máximo atómico, $\Psi_{C}$ el máximo en la componente
alcanzable $C$ y $\pi_{\beta,C}$ la estacionaria restringida. Para la cadena
reversible sobre $C$,
\begin{equation}
  \label{eq:cuatro-mecanismos}
  \Psi_{\mathrm{rel}}-\mathbb E_{\mu_t}[\Psi]\;\le\;
  \underbrace{\Psi_{\mathrm{rel}}-\Psi_{\mathrm{IP}}}_{\text{relajación}}
  +\underbrace{\Psi_{\mathrm{IP}}-\Psi_{C}}_{\text{topología}}
  +\underbrace{\frac{\log|C|}{\beta}}_{\text{temperatura}}
  +\underbrace{\mathrm{osc}_C(\Psi)\,\lVert\mu_t-\pi_{\beta,C}\rVert_{TV}}_{\text{mezcla}},
\end{equation}
y el último término decae al menos como $e^{-\gamma_Mt}$.
\end{proposicion}

\begin{proof}
Sumar y restar $\Psi_{\mathrm{IP}}$, $\Psi_C$ y $\mathbb E_\pi[\Psi]$; la cota
de temperatura $\Psi_C-\mathbb E_{\pi}[\Psi]\le\log|C|/\beta$ es la entropía
de Gibbs; $|\mathbb E_\pi-\mathbb E_\mu|\le\mathrm{osc}\cdot\lVert\mu-\pi\rVert_{TV}$.
\end{proof}

\begin{observacion}[Cuatro palancas que suelen mezclarse]
\label{obs:cuatro-palancas}
\Cref{pr:cuatro-mecanismos} refina \cref{th:descomposicion} en la capa
estocástica: mejorar el solver continuo ataca la relajación; ampliar el
soporte, la topología —y \cref{th:gap-monotono} dice que además acelera—;
subir $\beta$, la temperatura; esperar, la mezcla. Ni más explotación ni menos
ruido corrigen una componente equivocada: en los sesenta mundos difíciles del
banco, el soporte por pares fue desconectado con brecha cero en todos, el
soporte completo tuvo brecha positiva en todos, y no hubo una sola violación
de la monotonía $\gamma_{H_3}\ge\gamma_{H_2}$, $\gamma_{H_4}\ge\gamma_{H_3}$;
las medianas fueron $0$, $1{,}872$, $2{,}725$ y $2{,}523$ para $h=2,3,4$ y
soporte completo \cite{mayorga2026compendio}. La descomposición exacta sobre
ciento ochenta mundos da, en el estrato difícil, $3{,}78\%$ de relajación,
$4{,}32\%$ de topología y $0{,}43\%$ de dinámica, con error de identidad
inferior a $6\times10^{-17}$: la dinámica por pares agotó su término y quedó
exactamente sobre el piso topológico.
\end{observacion}
